\documentclass[11pt,a4paper]{article}
\usepackage[utf8]{inputenc}
\usepackage[french]{babel}
\usepackage[T1]{fontenc}
\usepackage{lmodern}
\usepackage[margin=1.5cm]{geometry}

\begin{document}

% En-tête
\begin{center}
    {\Huge \textbf{Emma Petit}} \\[0.2cm]
    {\Large \textbf{Architecte Big Data}} \\[0.1cm]
    emma.petit@email.com \ | \ 06 43 79 26 37
\end{center}

\vspace{0.3cm}

% Résumé / Profil
\section*{Profil Professionnel}
\noindent
Ingénieur IA passionné par le Traitement du Langage Naturel (NLP) et la Computer Vision. J'aide à transformer des données brutes en insights actionnables grâce à des architectures de réseaux de neurones complexes, tout en assurant leur mise en production via des API performantes. Convaincu par l'approche MLOps, je veille au maintien des performances des modèles tout au long de leur cycle de vie.

\vspace{0.3cm}

% Compétences
\section*{Compétences Techniques}
\noindent
\textbf{Deep Learning} $\bullet$ \textbf{TensorFlow} $\bullet$ \textbf{Spark} $\bullet$ \textbf{PyTorch} $\bullet$ \textbf{SQL} $\bullet$ \textbf{Pandas} $\bullet$ \textbf{Python} $\bullet$ \textbf{Machine Learning}

\vspace{0.3cm}

% Expériences
\section*{Expériences Professionnelles}
\noindent \textbf{Architecte Big Data / Confirmé} --- \textit{Criteo} \hfill \textbf{11/2021 à Aujourd'hui} \\
Collaboration avec les experts métiers pour définir les KPIs de succès et extraire les features les plus pertinentes des bases de données historiques. Industrialisation de l'extraction de données non structurées (OCR) sur des factures PDF via l'intégration d'API de Computer Vision. Conception d'un système de recommandation basé sur le Deep Learning, générant une hausse de 15\% du chiffre d'affaires sur le segment e-commerce. Entraînement et déploiement de modèles de Machine Learning (NLP) ayant permis d'automatiser 60\% des tâches manuelles de classification de documents. \\[0.3cm]
\noindent \textbf{Architecte Big Data} --- \textit{Dataiku} \hfill \textbf{11/2020 à 10/2021} \\
Mise en place d'une architecture de Feature Store centralisée, facilitant la réutilisation des variables entre les différentes équipes Data. Conception d'un modèle de prévision de la demande (Time Series) avec Prophet, réduisant les ruptures de stock de 30\% sur l'année. Optimisation des hyperparamètres de modèles Random Forest réduisant le taux de faux positifs de 22\% dans le système de détection de fraudes. Conception d'une architecture MLOps robuste sur AWS SageMaker pour garantir le suivi des performances et le réentraînement automatique des algorithmes. \\[0.3cm]
\noindent \textbf{Alternance — Assistant Architecte Big Data} --- \textit{Orange Business Services} \hfill \textbf{10/2019 à 04/2020} \\
Mise en place d'un pipeline Big Data avec Apache Spark pour traiter et nettoyer plus de 500 To de données par jour en temps réel. Déploiement de modèles prédictifs en production via FastAPI et Docker, avec intégration complète d'une boucle de feedback utilisateur. Création d'un algorithme de clustering non supervisé pour segmenter la base client en 5 catégories distinctes et optimiser les campagnes marketing. \\[0.3cm]


\vspace{0.3cm}

% Formations
\section*{Formation \& Diplômes}
\noindent \textbf{Diplôme d'Ingénieur en Mathématiques Appliquées} \hfill \textbf{2017 - 2020} \\
\textit{CentraleSupélec} \\[0.3cm]
\noindent \textbf{DUT en Mathématiques Appliquées} \hfill \textbf{2015 - 2017} \\
\textit{IUT de Bordeaux} \\[0.3cm]


\end{document}